\chapter{Generative Adversarial Networks}

\section{Introduction}

Generative Adversarial Networks (GANs) take a fundamentally different approach to generative modeling. Instead of explicitly modeling a likelihood, GANs learn to generate data by playing a game between two networks: a generator and a discriminator.

This chapter develops GANs from a game-theoretic perspective. We derive the minimax objective, analyze the role of the discriminator, and interpret GANs as implicitly minimizing a divergence between distributions. We also discuss training instabilities and common remedies.

\section{Adversarial Learning Setup}

GANs consist of two components:

\begin{itemize}
    \item \textbf{Generator} $G_\theta(z)$: maps noise $z \sim p(z)$ to data space,
    \item \textbf{Discriminator} $D_\phi(x)$: outputs a score indicating whether $x$ is real or generated.
\end{itemize}

The goal is:
\begin{itemize}
    \item $G_\theta$: produce realistic samples,
    \item $D_\phi$: distinguish real data from generated samples.
\end{itemize}

Thus training becomes a two-player game.

\section{Minimax Formulation}

The standard GAN objective is:
\[
\min_\theta \max_\phi \;
\mathbb{E}_{x \sim p_{\text{data}}}[\log D_\phi(x)]
+
\mathbb{E}_{z \sim p(z)}[\log(1 - D_\phi(G_\theta(z)))].
\]

\subsection{Interpretation}

\begin{itemize}
    \item the discriminator maximizes classification accuracy,
    \item the generator minimizes the discriminator's ability to distinguish samples.
\end{itemize}

This is a minimax optimization problem.

\section{What the Discriminator Learns}

For a fixed generator, the discriminator solves:
\[
\max_D \;
\mathbb{E}_{x \sim p_{\text{data}}}[\log D(x)]
+
\mathbb{E}_{x \sim p_G}[\log(1 - D(x))].
\]

\begin{proposition}[Optimal Discriminator]
The optimal discriminator is:
\[
D^*(x)
=
\frac{p_{\text{data}}(x)}{p_{\text{data}}(x) + p_G(x)}.
\]
\end{proposition}

\subsection{Interpretation}

\begin{itemize}
    \item $D^*(x)$ estimates the probability that $x$ is real,
    \item it compares densities of real and generated data.
\end{itemize}

Thus the discriminator implicitly learns a density ratio.

\section{Divergence Minimization View}

Substituting $D^*$ into the objective yields:
\[
\min_G \; 2 \cdot \mathrm{JS}(p_{\text{data}} \,\|\, p_G) - \log 4,
\]
where $\mathrm{JS}$ is the Jensen--Shannon divergence.

\subsection{Interpretation}

GAN training can be viewed as minimizing a divergence between:
\begin{itemize}
    \item the true data distribution,
    \item the generator distribution.
\end{itemize}

However, this divergence is implicit and not computed directly.

\section{Training Pathologies}

GAN training is notoriously unstable.

\subsection{Non-Convergence}

The minimax game may not converge:
\begin{itemize}
    \item oscillations between generator and discriminator,
    \item failure to reach equilibrium.
\end{itemize}

\subsection{Mode Collapse}

The generator produces limited diversity:
\begin{itemize}
    \item maps many $z$ values to similar outputs,
    \item ignores parts of the data distribution.
\end{itemize}

\subsection{Vanishing Gradients}

If the discriminator becomes too strong:
\[
D(x) \approx 1,\quad D(G(z)) \approx 0,
\]
then gradients for the generator become very small.

\section{Stabilization Tricks}

Several techniques improve GAN training:

\begin{itemize}
    \item \textbf{Non-saturating loss:} replace $\log(1-D(G(z)))$ with $-\log D(G(z))$,
    \item \textbf{Gradient penalties:} enforce smoothness of discriminator,
    \item \textbf{Label smoothing:} reduce overconfidence,
    \item \textbf{Architecture design:} use convolutional and residual structures,
    \item \textbf{Normalization:} stabilize training dynamics.
\end{itemize}

These do not fully solve instability but improve practical performance.

\section{Conditional GANs}

GANs can be extended to conditional settings:
\[
G(z, y), \quad D(x, y).
\]

\subsection{Objective}

\[
\min_G \max_D \;
\mathbb{E}_{x,y}[\log D(x,y)]
+
\mathbb{E}_{z,y}[\log(1 - D(G(z,y), y))].
\]

\subsection{Applications}

\begin{itemize}
    \item image-to-image translation,
    \item class-conditional generation,
    \item super-resolution.
\end{itemize}

Conditioning allows controlled generation.

\section{Worked Derivation: Optimal Discriminator}

We maximize:
\[
\mathcal{L}(D)
=
\int p_{\text{data}}(x)\log D(x)\,dx
+
\int p_G(x)\log(1-D(x))\,dx.
\]

For each $x$, maximize:
\[
p_{\text{data}}(x)\log D + p_G(x)\log(1-D).
\]

Setting derivative to zero:
\[
\frac{p_{\text{data}}(x)}{D(x)} - \frac{p_G(x)}{1-D(x)} = 0.
\]

Solving:
\[
D^*(x)=\frac{p_{\text{data}}(x)}{p_{\text{data}}(x)+p_G(x)}.
\]

\section{Wasserstein GAN (WGAN)}

Standard GANs minimize a Jensen--Shannon divergence, which can lead to vanishing gradients when the real and generated distributions have little overlap. Wasserstein GANs address this by replacing the divergence with the \emph{Wasserstein distance} (Earth Mover's distance), which provides more meaningful gradients.

\subsection{Wasserstein Distance (Intuition)}

The Wasserstein-1 distance between distributions $p$ and $q$ measures the minimum cost of transporting mass to transform $p$ into $q$.

\begin{itemize}
    \item robust even when supports do not overlap,
    \item provides informative gradients,
    \item better reflects distributional differences.
\end{itemize}

\subsection{WGAN Objective}

Using the Kantorovich--Rubinstein duality, the objective becomes:
\[
\min_G \max_{D \in \mathcal{D}}
\;
\mathbb{E}_{x \sim p_{\text{data}}}[D(x)]
-
\mathbb{E}_{z \sim p(z)}[D(G(z))],
\]
where $D$ (called a critic) is constrained to be 1-Lipschitz.

\subsection{Interpretation}

\begin{itemize}
    \item $D(x)$ assigns higher values to real samples,
    \item the difference estimates the Wasserstein distance,
    \item the generator minimizes this distance.
\end{itemize}

\section{Enforcing the Lipschitz Constraint}

The key requirement in WGAN is that the critic must be 1-Lipschitz.

\subsection{Weight Clipping (Original WGAN)}

The original method enforces this by clipping weights:
\[
w \leftarrow \mathrm{clip}(w, -c, c).
\]

\subsection{Limitations}

\begin{itemize}
    \item restricts model capacity,
    \item can lead to poor optimization,
    \item sensitive to choice of clipping threshold.
\end{itemize}

\section{WGAN-GP (Gradient Penalty)}

WGAN-GP improves stability by enforcing the Lipschitz constraint through a gradient penalty.

\subsection{Objective}

\[
\min_G \max_D
\;
\mathbb{E}_{x \sim p_{\text{data}}}[D(x)]
-
\mathbb{E}_{z \sim p(z)}[D(G(z))]
+
\lambda \, \mathbb{E}_{\hat{x}}
\left[
\left(\|\nabla_{\hat{x}} D(\hat{x})\|_2 - 1\right)^2
\right],
\]
where $\hat{x}$ are points interpolated between real and generated samples.

\subsection{Interpretation}

\begin{itemize}
    \item penalizes deviation from unit gradient norm,
    \item enforces approximate Lipschitz continuity,
    \item leads to more stable training.
\end{itemize}

\subsection{Advantages}

\begin{itemize}
    \item improved convergence,
    \item reduced mode collapse,
    \item more reliable gradients.
\end{itemize}

\section{Comparison with Standard GANs}

\begin{itemize}
    \item Standard GAN:
    \begin{itemize}
        \item minimizes JS divergence,
        \item prone to vanishing gradients,
        \item unstable training.
    \end{itemize}
    \item WGAN:
    \begin{itemize}
        \item minimizes Wasserstein distance,
        \item stable gradients,
        \item requires Lipschitz constraint.
    \end{itemize}
    \item WGAN-GP:
    \begin{itemize}
        \item improved enforcement via gradient penalty,
        \item more robust in practice.
    \end{itemize}
\end{itemize}

\section{Updated Perspective on GANs}

Wasserstein GANs highlight that the choice of divergence is crucial for training dynamics. By using a distance with better geometric properties, WGANs improve stability and interpretability.

This reinforces a broader principle:
\[
\text{good optimization} \;\Longleftrightarrow\; \text{well-behaved objective geometry}.
\]

Thus GAN variants can be understood as modifying the training objective to improve gradient behavior and convergence.

\section{A Unifying Perspective}

GANs replace likelihood-based learning with adversarial learning. The generator and discriminator form a game in which the generator improves by fooling the discriminator, while the discriminator improves by detecting generated samples.

This leads to:
\begin{itemize}
    \item powerful sample generation,
    \item implicit density modeling,
    \item challenging optimization dynamics.
\end{itemize}

GANs highlight a key idea: good generative models can be learned without explicitly computing likelihoods.

\section{Summary}

In this chapter we developed generative adversarial networks. We formulated GANs as a minimax game, derived the optimal discriminator, and interpreted training as divergence minimization. We discussed common training pathologies such as instability and mode collapse, along with practical remedies.

We also introduced conditional GANs for controlled generation. GANs provide a powerful alternative to likelihood-based generative modeling.